
\section*{7. Radicial Groups of Height 1}
\phantomsection
\label{sga3:VIIA:7}
\addcontentsline{toc}{section}{7. Radicial Groups of Height 1}

Let \(S\) be a scheme of characteristic
\(p>0\).\footnote{For the results of this section, one may also
consult [DG70], \S~II.7, nos.~3--4.}
We shall say that an \(\mathcal O_S\)-algebra \(\mathcal A\)
(respectively, an \(\mathcal O_S\)-\(p\)-Lie algebra
\(\mathcal L\)) is \emph{finite locally free} if its underlying
\(\mathcal O_S\)-module is locally free and of finite type. If
\(\mathcal L\) is a finite locally free
\(\mathcal O_S\)-\(p\)-Lie algebra, we know (cf.\ \SGARefLink{sga3:VIIA:5.5.2}{5.5.2}) that the
group \(S\)-functor
\[
  \operatorname{Spec}^{*}\mathcal U_p(\mathcal L)
\]
is represented by a finite locally free \(S\)-group scheme
\(\mathfrak G_p(\mathcal L)\). We shall see that this
\(S\)-group scheme solves a universal problem (\SGARefLink{sga3:VIIA:7.2}{7.2}), and we
shall characterize the \(S\)-group schemes of the form
\(\mathfrak G_p(\mathcal L)\) (\SGARefLink{sga3:VIIA:7.4}{7.4}).

\medskip
\noindent\textbf{Definition 7.0.}\footnote{The editors added
this definition (cf.\ [DG70], \S~II.4, \SGARefLink{sga3:VIIA:7.1}{7.1}), which will be used
in \SGARefLink{sga3:VIIA:7.2.1}{7.2.1}.}
\phantomsection
\label{sga3:VIIA:7.0}
\textit{Let \(H=\operatorname{Spec}\mathcal A\) be a finite
locally free \(S\)-group scheme. We say that \(H\) is
\emph{infinitesimal} if its unit section
\(\varepsilon_H:S\to H\) is a homeomorphism; equivalently, the
augmentation ideal of \(\mathcal A\) is locally nilpotent.}

\medskip
\noindent\textbf{7.1.}\footnote{The editors simplified the
original, taking into account the additions made in \SGARefLink{sga3:VIIA:5.5}{5.5}.}
\phantomsection
\label{sga3:VIIA:7.1}

Let \(\mathcal L\) be a finite locally free
\(\mathcal O_S\)-\(p\)-Lie algebra, and let
\(\mathfrak G_p(\mathcal L)\) be the affine \(S\)-group
\(\operatorname{Spec}\mathcal U_p(\mathcal L)\). By \SGARefLink{sga3:VIIA:5.5}{5.5},
\(\mathcal L\) is identified with
\(\operatorname{Lie}\mathfrak G_p(\mathcal L)\).

Now let \(G\) be a very good group \(S\)-functor satisfying
condition (F) of \SGARefLink{sga3:VIIA:6.3}{6.3}, and let
\[
  \varphi:\mathfrak G_p(\mathcal L)\longrightarrow G
\]
be a morphism of group \(S\)-functors. By \SGARefLink{sga3:VIIA:6.3}{6.3}, the morphism of
\(\mathcal O_S\)-Lie algebras
\[
  \operatorname{Lie}\varphi:
  \operatorname{Lie}\mathfrak G_p(\mathcal L)
  \longrightarrow
  \operatorname{Lie}G
\]
is compatible with the symbolic \(p\)-th power. If
\(\operatorname{Hom}_p(\mathcal L,\operatorname{Lie}G)\)
denotes the set of morphisms of \(\mathcal O_S\)-Lie algebras
compatible with the symbolic \(p\)-th power, we therefore have a
map
\[
  \operatorname{Lie}:
  \operatorname{Hom}_{S\text{-Gr.}}
    \bigl(\mathfrak G_p(\mathcal L),G\bigr)
  \longrightarrow
  \operatorname{Hom}_p
    \bigl(\mathcal L,\operatorname{Lie}G\bigr),
  \qquad
  \varphi\longmapsto\operatorname{Lie}\varphi.
\]

\medskip
\noindent\textbf{7.2. Theorem.}
\phantomsection
\label{sga3:VIIA:7.2}

\textit{If \(\mathcal L\) is a finite locally free
\(\mathcal O_S\)-\(p\)-Lie algebra, the map}
\[
  \operatorname{Hom}_{S\text{-gr.}}
    \bigl(\mathfrak G_p(\mathcal L),G\bigr)
  \longrightarrow
  \operatorname{Hom}_p
    \bigl(\mathcal L,\operatorname{Lie}G\bigr)
\]
\textit{is bijective in each of the following cases:}
\begin{enumerate}
  \item[\textit{(i)}]
  \textit{\(G\) is an \(S\)-group scheme;}
  \item[\textit{\SGAInlineRefTarget{sga3:VIIA:7.2:item:ii}(ii)}]
  \textit{\(G\) is of the form \(\operatorname{Aut}_S X\), where
  \(X\) is an \(S\)-scheme;}
  \item[\textit{(iii)}]
  \textit{\(G\) is of one of the forms
  \(\mathcal W(\mathcal F)\) or
  \(\operatorname{Aut}_{\mathcal O_S\text{-mod.}}
  \mathcal W(\mathcal F)\), where \(\mathcal F\) is a
  quasi-coherent \(\mathcal O_S\)-module.}
\end{enumerate}

The proof of the theorem rests on the following lemma.

\medskip
\noindent\textbf{Lemma.}
\phantomsection
\label{sga3:VIIA:7.2-frobenius-lemma}

\textit{If \(\mathcal L\) is a finite locally free
\(\mathcal O_S\)-\(p\)-Lie algebra, the \(S\)-group
\[
  G=\mathfrak G_p(\mathcal L)
\]
is annihilated by the Frobenius morphism
\(\operatorname{Fr}:G\to G^{(p)}\). In particular, \(G\) is
infinitesimal.}

Indeed, let \(\mathcal U\) be the restricted enveloping algebra
of \(\mathcal L\), let \(\mathcal A=\mathcal U^{*}\) be the
affine algebra of \(G\), and let
\(\mathcal I=\ker\varepsilon_{\mathcal A}\) be the augmentation
ideal of \(\mathcal A\).\footnote{The editors supplied details
in what follows. For another proof, see [DG70], \S~II.7, 3.9.}
We have
\begin{equation}\SGAInlineEquationTarget{sga3:viia:paragraph:7-2:equation:1}
  \mathcal A
  =
  \mathcal I\oplus\eta_{\mathcal A}(\mathcal O_S),
  \tag{1}
  \label{sga3:VIIA:7.2-lemma-decomposition-A}
\end{equation}
where \(\eta_{\mathcal A}\) denotes the unit section of
\(\mathcal A\). Since \(\varepsilon_{\mathcal A}\)
(respectively, \(\eta_{\mathcal A}\)) is the transpose of
\(\eta_{\mathcal U}\) (respectively,
\(\varepsilon_{\mathcal U}\)), this decomposition corresponds
by duality to
\begin{equation}\SGAInlineEquationTarget{sga3:viia:paragraph:7-2:equation:2}
  \mathcal U
  =
  \mathcal J\oplus\eta_{\mathcal U}(\mathcal O_S),
  \tag{2}
  \label{sga3:VIIA:7.2-lemma-decomposition-U}
\end{equation}
where \(\mathcal J\) is the augmentation ideal of
\(\mathcal U\); thus \(\mathcal J=\mathcal I^{*}\).

Let \(\pi\) denote the endomorphism
\(x\mapsto x^p\) of \(\mathcal O_S\). We must show that
\(\operatorname{Fr}:G\to G^{(p)}\) factors through the unit
section of \(G^{(p)}\). By \SGARefLink{sga3:VIIA:4.1.4}{4.1.4}~(c), this is equivalent to
saying that the morphism
\[
  \Phi:
  \mathcal I\otimes_{\pi}\mathcal O_S
  \longrightarrow
  \mathcal A,
  \qquad
  a\otimes_{\pi}x\longmapsto a^p x,
\]
is zero. Since \(\mathcal A\) is finite locally free over
\(\mathcal O_S\), it suffices to show that the transpose
\({}^{t}\Phi\) is zero.

Now \(\Phi\) is precisely the following composite:
\begin{center}
  \phantomsection
  \label{sga3:VIIA:7.2-frobenius-phi-composite}
  \(\displaystyle
\mathcal J\otimes_{\pi}\mathcal O_S
  \xrightarrow{\tau}
\mathcal A\otimes_{\pi}\mathcal O_S
  \xrightarrow{j(\mathcal A)}
\mathrm S^p\mathcal A
  \xrightarrow{b(\mathcal A)}
\mathcal A ,
\)

\end{center}
Here \(\tau\) is induced by the inclusion
\(\mathcal I\hookrightarrow\mathcal A\), while
\(b(\mathcal A)\) and \(j(\mathcal A)\) are defined as in
\SGARefLink{sga3:VIIA:4.3.3}{4.3.3}: \(b(\mathcal A)\) is induced by multiplication in
\(\mathcal A\), and \(j(\mathcal A)\) sends
\(a\otimes_{\pi}1\) to the image of
\(a\otimes\cdots\otimes a\) in
\(\mathrm S^p\mathcal A\). The
\(\mathcal O_S\)-dual of \(\mathrm S^p\mathcal A\) is the
submodule \(\Sigma^p\mathcal U\) of
\(\bigotimes^p\mathcal U\) consisting of the sections invariant
under the symmetric group of degree \(p\). Hence
\({}^{t}\Phi\) is the following composite:
\begin{center}
  \phantomsection
  \label{sga3:VIIA:7.2-transposed-phi-composite}
  \(\displaystyle
\mathcal U
  \xrightarrow{a(\mathcal U)}
\Sigma^p\mathcal U
  \xrightarrow{r(\mathcal U)}
\mathcal U\otimes_{\pi}\mathcal O_S
  \xrightarrow{q}
\mathcal J\otimes_{\pi}\mathcal O_S ,
\)

\end{center}
Here \(q\) is induced by the projection
\(\mathcal U\to\mathcal J\) with kernel
\(\eta_{\mathcal U}(\mathcal O_S)\),
\(a(\mathcal U)\) is induced by the comultiplication of
\(\mathcal U\), and \(r(\mathcal U)\) vanishes on symmetrized
tensors and sends a section
\(x\otimes\cdots\otimes x\) to
\(x\otimes_{\pi}1\) (cf.\ \SGARefLink{sga3:VIIA:4.3.3}{4.3.3}).

It is clear that
\({}^{t}\Phi\circ\eta_{\mathcal U}=0\); by (2), it remains to
show that \({}^{t}\Phi\) annihilates the augmentation ideal
\(\mathcal J\). Since \({}^{t}\Phi\) is an algebra morphism and
\(\mathcal J\) is generated by \(\mathcal L\), identified with
its image in \(\mathcal U\), it suffices to show that
\({}^{t}\Phi(x)=0\) for every section \(x\) of
\(\mathcal L\). But
\[
  -a(\mathcal U)(x)=(p-1)!\,a(\mathcal U)(x)
\]
is the symmetrization of
\(x\otimes1\otimes\cdots\otimes1\), so its image under
\(r(\mathcal U)\) is zero. This proves the first assertion of
the lemma.

The second follows. Every local section of \(\mathcal I\) has
zero \(p\)-th power, and \(\mathcal I\) is an
\(\mathcal O_S\)-module of finite type; hence \(\mathcal I\) is
locally nilpotent. Explicitly, if \(V\) is an affine open subset
of \(S\) such that
\[
  I=\Gamma(V,\mathcal I)
\]
is generated by \(r\) elements, then
\[
  I^{r(p-1)+1}=0.
\]
Thus \(G_{\mathrm{red}}=S_{\mathrm{red}}\), and the unit section
\(\varepsilon_G:S\to G\) is a homeomorphism.

\medskip
\noindent\textbf{7.2.1.}\footnote{The editors supplied details
in what follows.}
\phantomsection
\label{sga3:VIIA:7.2.1}

We first prove assertion \SGARefLink{sga3:vii:paragraph:5-1:equation:ii}{(ii)} of \SGARefLink{sga3:VIIA:7.2}{Theorem 7.2}. Let
\(\pi:X\to S\) be an \(S\)-scheme, and first consider an
arbitrary infinitesimal \(S\)-group \(H\).

Morphisms \(\varphi:H\to\operatorname{Aut}X\) correspond
bijectively to left actions
\(\mu:H\times X\to X\) of \(H\) on \(X\). For such an action,
if \(\varepsilon\) is the unit section of \(H\), the composite
\begin{center}
  \phantomsection
  \label{sga3:VIIA:7.2.1-action-unit-composite}
  \(\displaystyle
X \simeq S\times X
  \xrightarrow{\varepsilon\times X}
H\times X
  \xrightarrow{\mu}
X
\)

\end{center}
must be the identity. Since
\((H\times X)_{\mathrm{red}}\) is identified with
\(X_{\mathrm{red}}\), the morphism \(\mu\) must induce the
identity on the associated reduced schemes. In particular,
\(\mu\) induces an action of \(H\) on every open subset of
\(X\). Thus, for every open \(U\subset X\) affine over \(S\), we
obtain a morphism of unital associative algebras
\[
  \mathcal A(U)
  \longrightarrow
  \mathcal A(H)\otimes\mathcal A(U)
\]
which makes \(\mathcal A(U)\) a left
\(\mathcal A(H)\)-comodule, in such a way that the restriction
maps
\[
  \mathcal A(U)\longrightarrow\mathcal A(U'),
  \qquad U'\subset U,
\]
are comodule morphisms. Conversely, every datum of this kind
comes from a unique left action \(\mu:H\times X\to X\).

We also have the following lemma.

\medskip
\noindent\textbf{Lemma.}
\phantomsection
\label{sga3:VIIA:7.2.1-action-representation-lemma}

\textit{Let \(X=\operatorname{Spec}\mathcal C\) be an affine
\(S\)-scheme, let \(H=\operatorname{Spec}\mathcal A\) be an
infinitesimal \(S\)-group, and put
\[
  \mathcal U
  =
  \mathcal A^{*}
  =
  \operatorname{Hom}_{\mathcal O_S}
    (\mathcal A,\mathcal O_S).
\]
Left actions of \(H\) on \(X\) correspond bijectively to
representations of the algebra \(\mathcal U\) in the
\(\mathcal O_S\)-module \(\mathcal C\) such that}
\begin{align*}
  \textit{(a)}\quad&
  u(1_{\mathcal C})
  =
  \varepsilon(u)\,1_{\mathcal C},
  \\
  \textit{(b)}\quad&
  u(xy)
  =
  \sum_i v_i(x)w_i(y)
  \quad\textit{if}\quad
  \Delta u
  =
  \sum_i v_i\otimes w_i.
\end{align*}
\textit{Here \(u\) is an arbitrary section of \(\mathcal U\)
over an affine open subset \(V\) of \(S\), while \(x\) and \(y\)
are sections of \(\mathcal C\) over \(V\). The symbol
\(1_{\mathcal C}\) denotes the unit section of \(\mathcal C\),
and \(\varepsilon\) and \(\Delta\) denote the augmentation and
the diagonal morphism of \(\mathcal U\), respectively.}

Indeed, a left action \(\mu\) of \(H\) on \(X\) is defined by a
morphism of unital associative algebras
\[
  \lambda:\mathcal C\longrightarrow\mathcal A\otimes\mathcal C
\]
which makes \(\mathcal C\) a left
\(\mathcal A\)-comodule. Let \(\alpha\) be the composite
\begin{center}
  \phantomsection
  \label{sga3:VIIA:7.2.1-action-contraction-composite}
  \resizebox{0.98\linewidth}{!}{%
  \(\displaystyle
  \mathcal U\otimes_{\mathcal O_S}\mathcal C
    \xrightarrow{\mathcal U\otimes\lambda}
  \mathcal U\otimes_{\mathcal O_S}\mathcal A\otimes_{\mathcal O_S}\mathcal C
    \xrightarrow{\gamma\otimes\mathcal C}
  \mathcal O_S\otimes_{\mathcal O_S}\mathcal C
    \simeq \mathcal C
  \)%
}

\end{center}
where \(\gamma\) is the ``contraction'' from
\(\mathcal A^{*}\otimes_{\mathcal O_S}\mathcal A\) to
\(\mathcal O_S\). Since \(\mathcal A\) is finite locally free
over \(\mathcal O_S\), the map
\[
  \lambda
  \longmapsto
  (\gamma\otimes\mathcal C)(\mathcal U\otimes\lambda)
\]
is a bijection from
\[
  \operatorname{Hom}_{\mathcal O_S}
    (\mathcal C,\mathcal A\otimes\mathcal C)
  \quad\text{onto}\quad
  \operatorname{Hom}_{\mathcal O_S}
    (\mathcal U\otimes\mathcal C,\mathcal C).
\]
Moreover, by duality, the condition that \(\lambda\) define a
left \(\mathcal A\)-comodule structure (respectively, be a
morphism of unital associative algebras) is equivalent to the
condition that \(\alpha\) be a representation of
\(\mathcal U\) in \(\mathcal C\) (respectively, satisfy
conditions (a) and (b)). This proves the lemma.

For any representation of \(\mathcal U\) in the
\(\mathcal O_S\)-module \(\mathcal C\), it is also clear that
the sections \(u\) of \(\mathcal U\) satisfying conditions (a)
and (b) form a subalgebra of \(\mathcal U\).

In the particular case
\[
  H=\mathfrak G_p(\mathcal L)
\]
which concerns us, these conditions will therefore hold for
every section \(u\) of
\(\mathcal U=\mathcal U_p(\mathcal L)\) if they hold for the
sections \(u\) of \(\mathcal L\), identified with its image in
\(\mathcal U_p(\mathcal L)\).

If \(u\) is a section of \(\mathcal L\), conditions (a) and (b)
simply mean
\[
  u(1_{\mathcal C})=0,
  \qquad
  u(xy)=u(x)y+xu(y);
\]
that is, \(\alpha(u)\) is an \(\mathcal O_S\)-derivation of
\(\mathcal C=\mathcal A(X)\). Assertion \SGARefLink{sga3:VIIA:7.2:item:ii}{(ii)} of \SGARefLink{sga3:VIIA:7.2}{7.2} follows.
Indeed, every morphism
\[
  \varphi:\mathfrak G_p=\mathfrak G_p(\mathcal L)
  \longrightarrow\operatorname{Aut}X
\]
defines a morphism of \(p\)-Lie algebras
\(\operatorname{Lie}\varphi\) from
\(\mathcal L=\operatorname{Lie}\mathfrak G_p\) to
\[
  \pi_*
  \bigl(
    \operatorname{Der}_{\mathcal O_S}\mathcal O_X
  \bigr),
\]
and conversely every datum of this kind comes, by what precedes,
from a unique action
\(\mu:\mathfrak G_p(\mathcal L)\times X\to X\).

\medskip
\noindent\textbf{7.2.2.}
\phantomsection
\label{sga3:VIIA:7.2.2}

We now show how assertion \SGARefLink{sga3:vii:paragraph:5-1:equation:i}{(i)} of \SGARefLink{sga3:VIIA:7.2}{Theorem 7.2} follows from \SGARefLink{sga3:vii:paragraph:5-1:equation:ii}{(ii)}.
Let \(G\) be an \(S\)-group scheme. If \(T\) is an
\(S\)-scheme and \(x\in G(T)\), let
\(\ell_x^T\) (respectively, \(r_x^T\)) denote the left
(respectively, right) translation of \(G_T\) defined by \(x\).
The maps
\[
  \ell^T:x\longmapsto\ell_x^T
\]
therefore determine a homomorphism
\(\ell:G\to\operatorname{Aut}G\).

On the other hand, let \(f\) be a \(T\)-automorphism of \(G_T\),
and define
\[
  {}^x f=(r_x^T)^{-1}fr_x^T;
\]
thus, for every \(T'\to T\) and \(g\in G(T')\),
\[
  ({}^x f)(g)=f(gx)x^{-1}.
\]
In this way \(G\) acts on the left on the group \(S\)-functor
\(\operatorname{Aut}G\), and hence also on the functors
\[
  T\longmapsto
  \operatorname{Hom}_{T\text{-Gr.}}
    \bigl(\mathfrak G_p(\mathcal L_T),
          \operatorname{Aut}G_T\bigr)
\]
and
\[
  T\longmapsto
  \operatorname{Hom}_p
    \bigl(
      \mathcal L_T,
      \operatorname{Lie}(\operatorname{Aut}G_T/T)
    \bigr).
\]

The morphism
\(\ell^T:G_T\to\operatorname{Aut}G_T\) identifies \(G_T\) with
the group of automorphisms of the \(T\)-scheme \(G_T\) which
commute with right translations. Its derived morphism
\(\operatorname{Lie}(\ell^T)\) identifies
\(\operatorname{Lie}G_T\) with the \(p\)-Lie algebra of
\(\mathcal O_T\)-derivations of \(\mathcal O_{G_T}\) which
commute with right translations (cf.\ II, \SGARefLink{sga3:ii:scholium:4-11-1}{4.11.1}). They
therefore induce the commutative squares
\begin{center}
  \phantomsection
  \label{sga3:VIIA:7.2.2-lie-invariant-commutative-squares}
  \resizebox{0.98\linewidth}{!}{%
\begin{tikzcd}[ampersand replacement=\&,row sep=3.5em,column sep=4.5em]
\operatorname{Hom}_{T\text{-}\mathrm{Gr.}}
  \bigl(\mathfrak G_p(\mathcal L_T),G_T\bigr)
  \arrow[r,"\operatorname{Lie}"]
  \arrow[d,swap,"\ell_T"]
  \&
\operatorname{Hom}_p
  \bigl(\mathcal L_T,\operatorname{Lie}(G_T/T)\bigr)
  \arrow[d,"\operatorname{Lie}\ell_T"]
  \\
\operatorname{Hom}_{T\text{-}\mathrm{Gr.}}
  \bigl(\mathfrak G_p(\mathcal L_T),\underline{\operatorname{Aut}}\,G_T\bigr)
  \arrow[r,"\operatorname{Lie}"]
  \&
\operatorname{Hom}_p
  \bigl(\mathcal L_T,
    \operatorname{Lie}(\underline{\operatorname{Aut}}\,G_T/T)\bigr) .
\end{tikzcd}%
}

\end{center}
The images of the two vertical arrows are the subfunctors
formed by the invariants under the action of the \(S\)-group
\(G\). Since the lower horizontal arrow is invertible by \SGARefLink{sga3:VIIA:7.2.1}{7.2.1}
and is compatible with the action of \(G\), the upper
horizontal arrow is invertible as well. This proves \SGARefLink{sga3:VIIA:7.2}{7.2}~\SGARefLink{sga3:vii:paragraph:5-1:equation:i}{(i)}.

\medskip
\noindent\textbf{7.2.3.}
\phantomsection
\label{sga3:VIIA:7.2.3}

Consider the case in which
\(
  G=
  \operatorname{Aut}_{\mathcal O_S\text{-mod.}}
    \mathcal W(\mathcal F)
\).\footnote{In what follows, the editors expanded and
simplified the original, taking VIB, \SGARefLink{prop:sga3-VIB-11.6.1}{11.6.1} into account.}
Put
\[
  \mathcal U=\mathcal U_p(\mathcal L),
  \qquad
  \mathcal A=\mathcal U^{*},
  \qquad
  H=\mathfrak G_p(\mathcal L)=\operatorname{Spec}\mathcal A.
\]
Since \(H\) is affine over \(S\), VIB, \SGARefLink{prop:sga3-VIB-11.6.1}{11.6.1} shows that a
morphism of \(S\)-groups from \(H\) to
\(\operatorname{Aut}_{\mathcal O_S\text{-mod.}}
\mathcal W(\mathcal F)\) is the same as a right
\(\mathcal A\)-comodule structure
\[
  \mu:\mathcal F\longrightarrow\mathcal F\otimes\mathcal A.
\]
Since \(\mathcal A\) is finite locally free over
\(\mathcal O_S\), this is in turn equivalent to giving a
representation
\[
  \alpha:
  \mathcal U
  \longrightarrow
  \operatorname{End}_{\mathcal O_S}(\mathcal F)
\]
of \(\mathcal U\) in \(\mathcal F\). Finally, by the universal
property of \(\mathcal U=\mathcal U_p(\mathcal L)\), giving such
a morphism \(\alpha\) is equivalent to giving its restriction
\(\rho\) to \(\mathcal L\), identified with its image in
\(\mathcal U\), and \(\rho\) is a morphism of \(p\)-Lie
algebras from \(\mathcal L\) to
\(\operatorname{End}_{\mathcal O_S}(\mathcal F)\).

Finally, consider the case \(G=\mathcal W(\mathcal F)\), with
the preceding notation.\footnote{The editors added what
follows. The original said that ``the case of
\(\mathcal W(\mathcal F)\) is analogous.''}
To give a morphism of \(S\)-functors
\(\varphi:H\to\mathcal W(\mathcal F)\) is first of all
equivalent to giving an element \(\theta\) of
\(\Gamma(H,\mathcal F\otimes\mathcal O_H)\). Since \(H\) is
finite locally free over \(S\), we have
\[
  \Gamma(H,\mathcal F\otimes\mathcal O_H)
  =
  \Gamma(S,\mathcal F\otimes\mathcal A)
  =
  \operatorname{Hom}_{\mathcal O_S}(\mathcal U,\mathcal F).
\]
The condition that \(\varphi\) be a group morphism means that
\(\theta\), regarded as a morphism of \(\mathcal O_S\)-modules
\(\mathcal U\to\mathcal F\), vanishes on
\(1_{\mathcal U}\) and on \(\mathcal J^2\), where
\(\mathcal J\) is the augmentation ideal of \(\mathcal U\).
Consequently,
\begin{equation}
  \operatorname{Hom}_{S\text{-gr.}}
    \bigl(H,\mathcal W(\mathcal F)\bigr)
  =
  \operatorname{Hom}_{\mathcal O_S}
    \bigl(\mathcal J/\mathcal J^2,\mathcal F\bigr).
  \tag{1}
  \label{sga3:VIIA:7.2.3-hom-group-additive}
\end{equation}

Next consider the quasi-coherent sheaf
\([\mathcal L,\mathcal L]\), the image of the morphism
\[
  \mathcal L\otimes\mathcal L
  \longrightarrow
  \mathcal L,
  \qquad
  x\otimes y\longmapsto[x,y].
\]
For every affine open subset \(V\) of \(S\), we have
\[
  [\mathcal L,\mathcal L](V)
  =
  [\mathcal L(V),\mathcal L(V)].
\]
There is an exact sequence
\begin{center}
  \phantomsection
  \label{sga3:VIIA:7.2.3-augmentation-quotient-exact-sequence}
  \(\displaystyle
(\dagger)\qquad
0 \longrightarrow [\mathcal L,\mathcal L]
  \longrightarrow \mathcal L
  \xrightarrow{\pi}
\mathcal J/\mathcal J^2
  \longrightarrow 0 ,
\)

\end{center}
in which \(\pi\) is the composite of the inclusion
\(\mathcal L\hookrightarrow\mathcal J\) and the projection
\(\mathcal J\to\mathcal J/\mathcal J^2\).

Indeed, the question is local on \(S\), so we may suppose that
\(S\) is affine with ring \(R\), and that
\[
  L=\mathcal L(S)
\]
is free with basis \((x_1,\ldots,x_r)\). Identify \(L\) with
its image in \(U=U_p(L)\), and let \(K\) be the sub-\(R\)-module
of \(U\) which is the direct sum of \([L,L]\) and the submodule
with basis the monomials
\[
  x_1^{n_1}\cdots x_r^{n_r}
  \quad\text{such that}\quad
  n_1+\cdots+n_r\geq2.
\]
One checks that \(K\) is a two-sided ideal of \(U\). Since
\(K\) is contained in \(J^2\), where \(J\) is the augmentation
ideal of \(U\), and contains all the products \(x_ix_j\), which
generate \(J^2\), it follows that \(K=J^2\). Hence
\[
  J^2\cap L=[L,L],
\]
and the displayed sequence is exact.

By \SGARefLink{sga3:VIIA:6.4.2}{6.4.2},
\(\operatorname{Lie}\mathcal W(\mathcal F)\) is simply
\(\mathcal W(\mathcal F)\), with zero Lie bracket and zero
symbolic \(p\)-th power. From this and the preceding discussion,
we obtain
\begin{equation}
  \operatorname{Hom}_p(\mathcal L,\mathcal F)
  =
  \operatorname{Hom}_{\mathcal O_S}
    \bigl(\mathcal L/[\mathcal L,\mathcal L],\mathcal F\bigr)
  =
  \operatorname{Hom}_{\mathcal O_S}
    \bigl(\mathcal J/\mathcal J^2,\mathcal F\bigr).
  \tag{2}
  \label{sga3:VIIA:7.2.3-hom-p-additive}
\end{equation}
Together with \SGARefLink{sga3:VIIA:7.2.3-hom-group-additive}{(1)}, this completes the proof of \SGARefLink{sga3:VIIA:7.2}{Theorem 7.2}.

\medskip
\noindent\textbf{7.3. Lemma.}
\phantomsection
\label{sga3:VIIA:7.3}

\textit{If \(\mathcal L\) is a finite locally free
\(\mathcal O_S\)-\(p\)-Lie algebra, the morphism}
\[
  j_{\mathcal L}:
  \mathcal L
  \longrightarrow
  \operatorname{Lie}\mathfrak G_p(\mathcal L)
\]
\textit{of \SGARefLink{sga3:VIIA:5.5}{5.5} is invertible.}\footnote{To avoid changing the
numbering, the editors retained statement \SGARefLink{sga3:VIIA:7.3}{7.3}, although they
included it, together with its proof, in \SGARefLink{sga3:VIIA:5.5.1}{5.5.1}.}

For the proof, see \SGARefLink{sga3:VIIA:5.5.1}{5.5.1}.

\medskip
\noindent\textbf{7.4.}
\phantomsection
\label{sga3:VIIA:7.4}

To conclude this section, we give a characterization of the
\(S\)-group schemes of the form
\(\mathfrak G_p(\mathcal L)\), where \(\mathcal L\) is a finite
locally free \(\mathcal O_S\)-\(p\)-Lie algebra.

Let \(G\) be an \(S\)-group scheme, let \(\varepsilon_G\) be its
unit section, and let \(\mathcal I'\) be the kernel of the
morphism
\[
  \varepsilon_G^{-1}(\mathcal O_G)
  \longrightarrow
  \mathcal O_S
\]
corresponding to \(\varepsilon_G\). By \SGARefLink{sga3:VIIB:2.5:proposition}{2.5} and \SGARefLink{sga3:VIIA:1.3.1}{1.3.1}, the image
of \(\operatorname{Lie}(G/S)(S)\) in \(\mathrm U(G)\) is
identified with the morphisms of \(\mathcal O_S\)-modules from
\(\varepsilon_G^{-1}(\mathcal O_G)\) to \(\mathcal O_S\) which
vanish on the unit section of
\(\varepsilon_G^{-1}(\mathcal O_G)\) and on
\(\mathcal I'^{\,2}\).

We thereby recover the canonical isomorphism
\[
  \operatorname{Lie}(G/S)(S)
  \xrightarrow{\sim}
  \operatorname{Hom}_{\mathcal O_S}
    (\mathcal I'/\mathcal I'^{\,2},\mathcal O_S)
\]
of II, \SGARefLink{sga3:VIIA:3.3}{3.3} and \SGARefLink{sga3:ii:remark:4-11-4}{4.11.4}.\footnote{If \(G\) is affine over \(S\)
and \(\mathcal I\) denotes the augmentation ideal of
\(\mathcal A(G)\), then
\(\mathcal I'/\mathcal I'^{\,2}\) is identified with
\(\varepsilon_G^{*}(\mathcal I/\mathcal I^2)\); cf.\ loc. cit.}
As in the cited passage, put
\[
  \omega_{G/S}=\mathcal I'/\mathcal I'^{\,2}.
\]
Thus the sheaf \(\operatorname{Lie}(G/S)\) is identified with
\[
  \omega_{G/S}^{*}
  =
  \operatorname{Hom}_{\mathcal O_S}
    (\omega_{G/S},\mathcal O_S).
\]
Moreover,\footnote{The editors added the following sentence.}
if \(G=\mathfrak G_p(\mathcal L)\), where
\(\mathcal L\) is a finite locally free
\(\mathcal O_S\)-\(p\)-Lie algebra, then \SGARefLink{sga3:VIIA:5.5.3}{5.5.3} gives
\[
  \omega_{G/S}=\mathcal L^{*}.
\]

\medskip
\noindent\textbf{Theorem.}
\phantomsection
\label{sga3:VIIA:7.4-theorem}

\textit{If \(G\) is a group scheme over a scheme \(S\) of
characteristic \(p>0\), the following assertions are
equivalent:}
\begin{enumerate}
  \item[\textit{\SGAInlineRefTarget{sga3:VIIA:7.4-theorem:item:i}(i)}]
  \textit{There exists a finite locally free
  \(\mathcal O_S\)-\(p\)-Lie algebra \(\mathcal L\) such that
  \(G\simeq\mathfrak G_p(\mathcal L)\).}

  \item[\textit{(i\(^{\prime}\))}]
  \textit{The \(\mathcal O_S\)-\(p\)-Lie algebra
  \(\operatorname{Lie}(G)\) is finite locally free, and
  \(G\simeq\mathfrak G_p(\operatorname{Lie}(G))\).}

  \item[\textit{(ii)}]
  \textit{\(G\) is affine over \(S\), the
  \(\mathcal O_S\)-module \(\omega_{G/S}\) is locally free of
  finite type, and the affine algebra of \(G\) is locally
  isomorphic to the quotient of the symmetric algebra
  \(\mathrm S_{\mathcal O_S}(\omega_{G/S})\) by the ideal
  generated by the \(p\)-th powers of the sections of
  \(\omega_{G/S}\).}

  \item[\textit{\SGAInlineRefTarget{sga3:VIIA:7.4-theorem:item:iii}(iii)}]
  \textit{\(G\) is locally of finite presentation over \(S\),
  is of height \(\leq1\), and \(\omega_{G/S}\) is locally free.}

  \item[\textit{(iii\(^{\prime}\))}]
  \textit{\(G\) is locally of finite type over \(S\), is of
  height \(\leq1\), and \(\omega_{G/S}\) is locally free.}

  \item[\textit{\SGAInlineRefTarget{sga3:VIIA:7.4-theorem:item:iv}(iv)}]
  \textit{\(G\) is locally of finite presentation and flat over
  \(S\), and is of height \(\leq1\).}\footnote{The editors added
  assertion (i\(^{\prime}\)), implicit in the original, and
  assertions (iii\(^{\prime}\)) and \SGARefLink{sga3:VIIA:7.4-theorem:item:iv}{(iv)}, pointed out by
  O.~Gabber. Assertion \SGARefLink{sga3:VIIA:7.4-theorem:item:iv}{(iv)} takes up a footnote in the original,
  which stated: ``The condition on \(\omega_{G/S}\) is in fact
  unnecessary, as is easily seen by reducing to the case where
  \(S\) is local with residue field \(k\), and applying the
  theorem to \(G_k\).'' As Gabber pointed out, this is incorrect
  without a flatness hypothesis. Indeed, let \(A\) be an
  Artinian local ring of characteristic \(p>0\), let \(J\) be a
  proper nonzero ideal of \(A\), and let \(H\) be the subgroup
  \[
    \operatorname{Spec}A[x]/(x^p,Jx)
  \]
  of \(\alpha_{p,A}\); that is, for every \(A\)-algebra \(R\),
  \[
    H(R)=\{x\in R\mid x^p=0\ \text{and}\ Jx=0\}.
  \]
  Then \(H\) is not flat over \(A\), and hence is not of the
  form \(\mathfrak G_p(\mathcal L)\) with \(\mathcal L\) a
  finite-rank free \(p\)-Lie algebra over \(A\).}
\end{enumerate}

\medskip
\noindent\textbf{7.4.1.}
\phantomsection
\label{sga3:VIIA:7.4.1}

The equivalence \SGARefLink{sga3:vii:paragraph:5-1:equation:i}{(i)} \(\Leftrightarrow\)
(i\(^{\prime}\)) follows from \SGARefLink{sga3:VIIA:5.5.3}{5.5.3}~\SGARefLink{sga3:vii:paragraph:5-1:equation:i}{(i)}, and the implications
\SGARefLink{sga3:VIIA:5.5.3:item:ii}{(ii)} \(\Rightarrow\) \SGARefLink{sga3:VIIA:7.4-theorem:item:iii}{(iii)} \(\Rightarrow\)
(iii\(^{\prime}\)) are clear. We also have
\SGARefLink{sga3:VIIA:5.5.3:item:i}{(i)} \(\Rightarrow\) \SGARefLink{sga3:VIIA:7.4-theorem:item:iv}{(iv)}, since
\(\mathfrak G_p(\mathcal L)\) is finite locally free and of
height \(\leq1\), by \SGARefLink{sga3:VIIA:5.5.2}{5.5.2} and the lemma in \SGARefLink{sga3:VIIA:7.2}{7.2}. We show that
\SGARefLink{sga3:vii:paragraph:5-1:equation:i}{(i)} implies \SGARefLink{sga3:vii:paragraph:5-1:equation:ii}{(ii)}.

Let \(\mathcal I\) be the augmentation ideal of
\[
  \mathcal A=\mathcal U_p(\mathcal L)^{*}.
\]
We saw in \SGARefLink{sga3:VIIA:5.5.3}{5.5.3}~\SGARefLink{sga3:vii:paragraph:5-1:equation:ii}{(ii)} that
\(\omega_{G/S}=\mathcal I/\mathcal I^2\) is identified with
\(\mathcal L^{*}\), and hence is finite locally free.

Now suppose that \(S\) is affine. There is a section
\[
  \sigma:\omega_{G/S}\longrightarrow\mathcal I
\]
of the projection
\(\mathcal I\to\mathcal I/\mathcal I^2\). It induces an algebra
morphism
\[
  \sigma':
  \mathrm S_{\mathcal O_S}(\omega_{G/S})
  \longrightarrow
  \mathcal A,
\]
and, by the lemma in \SGARefLink{sga3:VIIA:7.2}{7.2}, \(\sigma'\) factors through a morphism
\[
  \varphi:
  \mathrm S_{\mathcal O_S}(\omega_{G/S})/\mathcal K
  \longrightarrow
  \mathcal A,
\]
where \(\mathcal K\) denotes the ideal generated by the
\(p\)-th powers of sections of \(\omega_{G/S}\). Filter
\(\mathcal A\) by powers of \(\mathcal I\), and filter
\(\mathrm S_{\mathcal O_S}(\omega_{G/S})/\mathcal K\) by
powers of the ideal generated by \(\omega_{G/S}\). The morphism
\(\varphi\) plainly induces an epimorphism on the associated
graded modules. Hence \(\varphi\) is an epimorphism between
locally free \(\mathcal O_S\)-modules of the same rank
(cf.\ \SGARefLink{sga3:VIIA:5.3.3}{5.3.3}), and is therefore an isomorphism. This proves that
\SGARefLink{sga3:vii:paragraph:5-1:equation:i}{(i)} implies \SGARefLink{sga3:vii:paragraph:5-1:equation:ii}{(ii)}.

\medskip
\noindent\textbf{7.4.2.}
\phantomsection
\label{sga3:VIIA:7.4.2}

Suppose now that \(G\) is of height \(\leq1\) and locally of
finite presentation over \(S\).\footnote{The editors supplied
details and simplified the original in what follows.}
The Frobenius morphism
\(\operatorname{Fr}:G\to G^{(p)}\) is integral and factors
through the unit section of \(G^{(p)}\); consequently, \(G\) is
integral, and hence affine, over \(S\). Put
\(\mathcal A=\mathcal A(G)\). Since \(G\) is locally of finite
presentation over \(S\), it follows that \(G\) is finite and of
finite presentation over \(S\), and hence that
\(\mathcal A(G)\) is an \(\mathcal O_S\)-module of finite
presentation (cf.\ EGA IV\(_1\), 1.4.7).

Let \(\mathcal I\) be the augmentation ideal of
\(\mathcal A\). Since
\[
  \mathcal A
  =
  \eta_{\mathcal A}(\mathcal O_S)\oplus\mathcal I,
\]
where \(\eta_{\mathcal A}\) is the unit section of
\(\mathcal A\), the module \(\mathcal I\) is of finite
presentation over \(\mathcal O_S\), as is
\(\omega_G=\mathcal I/\mathcal I^2\). If instead \(G\) is
assumed to be of height \(\leq1\) and locally of finite type
over \(S\), the same argument shows that
\(\mathcal A(G)\), \(\mathcal I\), and
\(\omega_G=\mathcal I/\mathcal I^2\) are
\(\mathcal O_S\)-modules of finite type.

Thus, under hypothesis (iii\(^{\prime}\)),
\(\omega_{G/S}\) is finite locally free over
\(\mathcal O_S\), as is
\[
  \mathcal L
  =
  \operatorname{Lie}(G/S)
  =
  \omega_{G/S}^{*}.
\]
Put
\[
  \mathcal B=\mathcal U_p(\mathcal L)^{*},
  \qquad
  H=\mathfrak G_p(\mathcal L)=\operatorname{Spec}\mathcal B.
\]
By \SGARefLink{sga3:VIIA:7.2}{Theorem 7.2}, the identity map of \(\mathcal L\) corresponds
to a group morphism from
\(H=\mathfrak G_p(\mathcal L)\) to \(G\), and hence to a
morphism of \(\mathcal O_S\)-algebras
\[
  \theta:\mathcal A\longrightarrow\mathcal B.
\]
By definition, \(\theta\) induces an isomorphism from
\(\omega_{G/S}\) to \(\omega_{H/S}\); we must show that
\(\theta\) itself is an isomorphism.

We may restrict to the case in which \(S\) is affine. There is
then a section \(\tau\) of the projection
\(\mathcal I\to\omega_{G/S}\), inducing an algebra morphism
\[
  \tau':
  \mathrm S_{\mathcal O_S}(\omega_{G/S})
  \longrightarrow
  \mathcal A.
\]
Every local section of \(\mathcal I\) has zero \(p\)-th power,
because \(\operatorname{Fr}:G\to G^{(p)}\) factors through the
unit section of \(G^{(p)}\). Thus \(\tau'\) induces a morphism
of \(\mathcal O_S\)-algebras \(\psi\) fitting into the
commutative diagram
\begin{center}
  \phantomsection
  \label{sga3:VIIA:7.4.2-height-one-characterization-triangle}
  \begin{tikzcd}[row sep=3.5em,column sep=4.6em]
\mathrm S_{\mathcal O_S}(\omega_{G/S})/\mathcal K
  \arrow[r,"\psi"]
  \arrow[dr,swap,"\varphi"]
  & \mathcal A \arrow[d,"\theta"] \\
& \mathcal B
\end{tikzcd}

\end{center}
where \(\mathcal K\) is the ideal generated by the \(p\)-th
powers of sections of \(\omega_{G/S}\). As in \SGARefLink{sga3:VIIA:7.4.1}{7.4.1},
\(\psi\) is an epimorphism of \(\mathcal O_S\)-modules. On the
other hand, \SGARefLink{sga3:VIIA:7.4.1}{7.4.1} shows that
\(\varphi=\theta\circ\psi\) is an isomorphism. Hence
\(\theta\) is an isomorphism as well. This proves
(iii\(^{\prime}\)) \(\Rightarrow\) \SGARefLink{sga3:VIIA:7.4-theorem:item:i}{(i)}.

\medskip
\noindent\textbf{7.4.3.}\footnote{The editors added this
paragraph to prove \SGARefLink{sga3:VIIA:7.4-theorem:item:iv}{(iv)} \(\Rightarrow\) \SGARefLink{sga3:VIIA:7.4-theorem:item:iii}{(iii)}; cf.\ editor's
note~\SGARefLink{sga3:VIIB:nde:71}{(71)}.}
\phantomsection
\label{sga3:VIIA:7.4.3}

We finally prove that \SGARefLink{sga3:VIIA:7.4-theorem:item:iv}{(iv)} implies \SGARefLink{sga3:VIIA:7.4-theorem:item:iii}{(iii)}. It suffices to show
that \(\omega_{G/S}\) is locally free, so we may suppose that
\(S\) is affine with ring \(R\). As observed at the beginning
of \SGARefLink{sga3:VIIA:7.4.2}{7.4.2}, hypothesis \SGARefLink{sga3:VIIA:7.4-theorem:item:iv}{(iv)} then implies that
\(G=\operatorname{Spec}A\), where \(A\) is an \(R\)-algebra
which is an \(R\)-module of finite presentation, as is
\[
  \omega_{G/S}=I/I^2,
\]
where \(I\) is the augmentation ideal of \(A\). Since \(G\) is
also assumed flat over \(S\), the algebra \(A\) is a finite
locally free \(R\)-module (cf.\ [BAC] II, \S~\SGARefLink{sga3:VIIB:5.2:theorem}{5.2}, Theorem~1 and
Corollary~2). By the cited result, to show that
\(\omega_{G/S}\) is locally free of finite rank, it suffices to
show that
\[
  (\omega_{G/S})_{\mathfrak m}
\]
is flat for every maximal ideal \(\mathfrak m\) of \(R\).
Thus we may suppose that \(R\) is local and that \(A\) is free
of rank \(n+1\); then \(I\) is free of rank \(n\). Let
\(\mathfrak m\) be the maximal ideal of \(R\), and put
\[
  k=R/\mathfrak m.
\]

Let \(I_k\) be the augmentation ideal of \(A_k\), and let \(r\)
be the dimension of
\[
  \omega_{G_k/k}=I_k/I_k^2.
\]
Choose a basis \((e_1,\ldots,e_n)\) of \(I_k\) such that
\((e_{r+1},\ldots,e_n)\) is a basis of \(I_k^2\), and choose
elements \(x_1,\ldots,x_n\) of \(I\) lifting the \(e_i\).
Nakayama's lemma shows that \((x_1,\ldots,x_n)\) is a basis of
\(I\) over \(R\).

Let \(N\) be the sub-\(R\)-module of \(I\) with basis
\((x_1,\ldots,x_r)\), and let \(B\) be the quotient of the
symmetric algebra of \(N\) by the ideal generated by the
elements \(x^p\), for \(x\in N\). Since every element of \(I\)
has zero \(p\)-th power, we obtain a morphism of \(R\)-algebras
\[
  \psi:B\longrightarrow A.
\]
By \SGARefLink{sga3:VIIA:7.4.2}{7.4.2}, \(\psi\otimes k\) is an isomorphism. It follows that
\(\operatorname{Coker}\psi=0\), and, if
\(K=\ker\psi\), the morphism
\[
  \tau:K\otimes k\longrightarrow B\otimes k
\]
is zero. But \(\psi\) is surjective and \(A\) is flat over
\(R\), so \(\tau\) is also injective. Hence
\(K\otimes k=0\).

Since \(A\) is an \(R\)-module of finite presentation, \(K\) is
an \(R\)-module of finite type (cf.\ [BAC] I, \S~\SGARefLink{sga3:VIIB:2.8}{2.8},
Lemma~9). Nakayama's lemma now gives \(K=0\). Thus \(\psi\) is
an isomorphism of \(R\)-algebras. Since \(\psi^{-1}(I)\)
contains the augmentation ideal \(J\) of \(B\), it follows that
\(\psi^{-1}(I)=J\). Therefore \(\psi^{-1}\) induces an
isomorphism of \(R\)-modules
\[
  I/I^2
  \xrightarrow{\sim}
  J/J^2=N.
\]
This proves that \(\omega_{G/S}\) is finite locally free, and
hence proves \SGARefLink{sga3:VIIA:7.4-theorem:item:iv}{(iv)} \(\Rightarrow\) \SGARefLink{sga3:VIIA:7.4-theorem:item:iii}{(iii)}. The proof of
\SGARefLink{sga3:VIIA:7.4}{Theorem 7.4} is complete.

\medskip
\noindent\textbf{Remark 7.5.}\footnote{The editors added this
remark.}
\phantomsection
\label{sga3:VIIA:7.5}

It follows immediately from \SGARefLink{sga3:VIIA:7.2}{Theorems 7.2} and \SGARefLink{sga3:VIIA:7.4}{7.4} that the
functors
\[
  G\longmapsto\operatorname{Lie}(G),
  \qquad
  \mathcal L\longmapsto\mathfrak G_p(\mathcal L)
\]
induce quasi-inverse equivalences between the category of
\(S\)-groups which are locally of finite presentation, flat,
and of height \(\leq1\), and the full subcategory of the
category of \(\mathcal O_S\)-\(p\)-Lie algebras consisting of
the finite locally free \(\mathcal O_S\)-\(p\)-Lie algebras.

